\documentclass[../main.tex]{subfiles}

\begin{document}

    \section{Why are function calls expensive?}\label{sec:why-are-function-calls-expensive}

    Python is an interpreted language (the interpreter standardly being CPython), which at runtime must:
    \begin{itemize}
        \item Parse the source code into an abstract syntax tree (AST).
        \item Compile this into bytecode - these are a set of instructions for the computer to follow.
        \item Execute that bytecode in an evaluation loop.
    \end{itemize}\\
    The overhead that function calls bring are mostly in the execution stage of runtime.
    So, let's look at what happens when you call a simple function \verb|y = f(x)| during execution:
    \begin{itemize}
        \item Step 1: Evaluate the callable object \verb|f| (our function object).
        \begin{itemize}
            \item CPython looks up the name \verb|f| in each of the scopes of the program:
            \begin{itemize}
                \item \textbf{locals} - the variables within the current function, whatever that may be.
                \item \textbf{enclosing scopes} - if the current function is nested (written inside another function with local variables).
                \item \textbf{globals} - functions within the top scope (for example, constants defined at the top of the file).
                \item \textbf{builtins} - functions/global variables like \verb|range|, \verb|max|, \verb|min|, etc.
            \end{itemize}\\
            \item The result from this lookup \textbf{must} be a callable object since we are actively calling \verb|f| with \verb|()| parenthesis.
            \item This process is also not instant and add a cost to our function call already.
        \end{itemize}\\
        \item Step 2: The arguments passed, in this case \verb|x|, must then be evaluated.
        \begin{itemize}
            \item The expression within the parenthesis must is evaluated from left to right.
            \item This could mean more name lookups (like above in Step 1), attribute lookups, or indexing.
            \item The resulting argument values are then pushed onto the value stack.
            \item This is an expensive part of the execution (evaluation and pushing to the stack).
        \end{itemize}\\
        \item Step 3: Building a call signature.
        \begin{itemize}
            \item \verb|f(x)| - creates a simple array of arguments.
            \item \verb|f(x, y=3)| - builds a map of the keyword names (kwargs) to values (args).
            \item \verb|f(*arg, **kwargs)| - unpacks mappings.
        \end{itemize}\\
        \item Step 4: Actually call the function - this uses a generic call structure.
        \begin{itemize}
            \item Checks if the object (\verb|f|) is callable (if it is a function).
            \item Finds the appropriate entry (where to start execution from) to the function.
            \item Fortunately, builtins are specifically made to be fast (using special routing) and use C-based methods, so these are less expensive to call.
        \end{itemize}
        \item Step 5: Creates a new, rather expensive, stack frame. This is an environment for the function to run within, containing:
        \begin{itemize}
            \item An instruction pointer - this tracks your position in the bytecode.
            \item Local variables - this is for fast lookups described partly in Step 2.
            \item Global/builtin references - so the function isn't limited to the local scope.
            \item An evaluation stack.
            \item Exception handling information.
        \end{itemize}
        \item Step 6: Binds arguments to parameters.
        \begin{itemize}
            \item CPython maps these onto the function's parameters:
            \begin{itemize}
                \item Positional parameters.
                \item Keyword-only parameters.
                \item Default values.
                \item args and kwargs.
            \end{itemize}
            \item While also checking for errors:
            \begin{itemize}
                \item Missing required arguments.
                \item Unexpected arguments/keywords.
                \item Other general checks.
            \end{itemize}
        \end{itemize}
        \item Step 7: Finally start to execute the bytecode.
        \begin{itemize}
            \item Initialisation of local variables.
            \item Execution begins at the address of the entry point for \verb|f|.
            \item CPython starts a bytecode evaluation loop for this new frame:
            \begin{itemize}
                \item Fetch the next bytecode instruction.
                \item Decode instruction into an opcode.
                \item Dispatch (jump) to a C code implementation of that opcode.
                \item Manipulate the value stack to reflect result of operation.
                \item Increment the instruction pointer and repeat.
            \end{itemize}
        \end{itemize}
        \item Step 8: Create return value and teardown frame.
        \begin{itemize}
            \item Produces a return value object if necessary.
            \item Removes the frame from the call stack.
            \item Restores calling frame back to the current frame.
            \item Pushes the return value onto the caller's evaluation stack.
        \end{itemize}
        \item Step 9: Assign result to \verb|y| - updates the local scope binding (of the calling scope) of \verb|y|.
    \end{itemize}\\
    Comparing with the steps for inline arithmetic, for example with \verb|y = x*x + 2|:
    \begin{itemize}
        \item Fetch next bytecode instruction.
        \item Lookup \verb|x|, which is found within the \textbf{locals} array (this is the first lookup array, so extremely fast), and push the value of \verb|x| onto the value stack.
        \\ \verb|stack = [x]|
        \item Lookup \verb|x| a second time and push onto the stack.
        \\ \verb|stack = [x, x]|
        \item Multiply \verb|x * x|.
        \begin{itemize}
            \item Pop top two items from stack (\verb|x| and \verb|x|).
            \\ \verb|stack = []|
            \item Determine the correct multiplication operation for the runtime type of the variables (\verb|int|, \verb|float|, etc.).
            \item Dispatch to a C code implementation (multiplication in this case).
            \item Push result onto the stack.
            \\ \verb|stack = [x_squared]|
        \end{itemize}
        \item Load \verb|2| and push onto the stack.
        \\ \verb|stack = [x_squared, 2]|
        \item Add \verb|x_squared + 2|.
        \begin{itemize}
            \item Pop top two items from stack (\verb|x_squared| and \verb|2|).
            \\ \verb|stack = []|
            \item Determine the correct addition operation for the runtime type of the variables (\verb|int|, \verb|float|, etc.).
            \item Dispatch to a C code implementation (addition in this case).
            \item Push result onto the stack.
            \\ \verb|stack = [result]|
        \end{itemize}
        \item Store result into \verb|y|.
        \begin{itemize}
            \item Pop top item from stack (\verb|result|).
            \\ \verb|stack = []|
            \item Store \verb|result| in \verb|y|.
        \end{itemize}
    \end{itemize}

    As you can see, inline arithmetic is much cheaper for the CPU as it doesn't require argument evaluation,
    building call signatures, stack frame creation, binding of arguments, error handling, or creation of return values.
    These are all relatively expensive operations and should be considered when choosing methods or algorithms when large amounts
    of data or iterations are required.

\end{document}